\subsection{Boundary Regularity}\label{sec:elliptic-boundary-regularity}

In the previous section we showed that if the data and coefficients are sufficiently smooth, then weak solutions of elliptic problems are ``as smooth as one could expect\textquotesingle\textquotesingle\ in the interior of the domain on which the problem is posed. In the following example we see that a solution is not necessarily smooth up to the boundary of the domain on which the problem is posed if that boundary is not sufficiently smooth. (The reader should also note that this is not some sort of weird ``cooked up\textquotesingle\textquotesingle\ counterexample; the domain in question simply has a corner.)

\begin{Example}
\label{ex:nonsmooth-sector-boundary-regularity}
Let $\bm{i}$ and $\bm{j}$ give an orthonormal basis for $\mathbb{R}^2$, and let $(r,\theta)$ be the standard polar coordinates for $\mathbb{R}^2$ defined by the map
\begin{equation}\label{eq:corner-example-polar-coordinate-map}
\widehat{\bm{x}}(r,\theta):=r\bm{e}_1(\theta),
\end{equation}
where
\begin{equation}\label{eq:corner-example-first-polar-basis-vector}
\bm{e}_1(\theta):=\cos\theta\bm{i}+\sin\theta\bm{j},
\end{equation}

\begin{equation}\label{eq:corner-example-second-polar-basis-vector}
\bm{e}_2(\theta):=-\sin\theta\bm{i}+\cos\theta\bm{j}.
\end{equation}
Recall that for a real-valued function $f(r,\theta)$ we can calculate the gradient and Laplacian as follows:
\begin{equation}\label{eq:corner-example-polar-gradient-formula}
\nabla f=f_r\bm{e}_1+\frac{1}{r}f_\theta\bm{e}_2,
\end{equation}

\begin{equation}\label{eq:corner-example-polar-laplacian-formula}
\Delta f=\frac{1}{r}\frac{\partial}{\partial r}\left(r\frac{\partial f}{\partial r}\right)+\frac{1}{r^2}\frac{\partial^2f}{\partial\theta^2}.
\end{equation}
In addition, the Hessian or second gradient matrix is given by
\begin{equation}\label{eq:corner-example-polar-hessian-matrix}
\begin{pmatrix}
f_{rr}&\dfrac{1}{r}f_{r\theta}-\dfrac{1}{r^2}f_\theta\\
\dfrac{1}{r}f_{r\theta}-\dfrac{1}{r^2}f_\theta&\dfrac{1}{r}f_r+\dfrac{1}{r^2}f_{\theta\theta}
\end{pmatrix}.
\end{equation}
We now consider the following problem for Laplace\textquotesingle s equation with nonhomogeneous boundary conditions. Let $0<\beta<2\pi$ and define
\begin{equation}\label{eq:corner-example-planar-sector-domain}
\Omega_\beta:=\{\bm{x}\in\mathbb{R}^2\mid\bm{x}=\widehat{\bm{x}}(r,\theta),\quad0<r<1,\quad0<\theta<\beta\}.
\end{equation}
We seek $u:\Omega_\beta\to\mathbb{R}$ satisfying
\begin{equation}\label{eq:corner-example-harmonic-sector-equation}
\Delta u=0\quad\text{in }\Omega_\beta,
\end{equation}
and the boundary conditions\[\begin{aligned}
u(r,0)&=0,&&0<r<1\\
u(r,\beta)&=0,&&0<r<1\\
u(1,\theta)&=\sin\frac{\pi\theta}{\beta},&&0<\theta<\beta.
\end{aligned}\]Using separation of variables, we can find the solution of the problem to be
\begin{equation}\label{eq:corner-example-explicit-harmonic-sector-solution}
u(r,\theta)=r^{\pi/\beta}\sin\frac{\pi\theta}{\beta}.
\end{equation}
This function (as the results of the previous section assure us) is in $C^\infty(\Omega^{\prime})$ for any $\Omega^{\prime}$ compactly contained in $\Omega_\beta$. However, note that if $\beta>\pi$, then our solution decays at the origin like $r^\alpha$ with $\alpha<1$. To see the impact of this we calculate the gradient of $u$
\begin{equation}\label{eq:corner-example-explicit-solution-gradient}
\nabla u=\frac{\pi}{\beta}r^{\pi/\beta-1}\left(\sin\frac{\pi\theta}{\beta}\bm{e}_1(\theta)+\cos\frac{\pi\theta}{\beta}\bm{e}_2(\theta)\right)
\end{equation}
and its $L^2(\Omega_\beta)$ norm
\begin{equation}\label{eq:corner-example-finite-gradient-l2-norm}
\int_{\Omega_\beta}|\nabla u|^2\,\mathrm{d}\bm{x}=C\int_0^\beta\int_0^1r^{2(\pi/\beta-1)}r\,\mathrm{d}r\,\mathrm{d}\theta<\infty.
\end{equation}
We see that (as our basic existence theorem for weak solutions implies) our solution $u$ is in $H^1(\Omega_\beta)$. However, by calculating the second gradient and computing its norm, we see that if $\beta>\pi$, then $u$ is not in $H^2(\Omega_\beta)$. Thus, despite the fact that we have all of the interior regularity guaranteed by the results of the previous section, we do not have regularity up to the boundary. The culprit here is the lack of smoothness of the boundary.
\end{Example}
